\documentclass[9pt,twocolumn,a4paper]{extarticle}

% ELECO 2026 yazım düzeni: A4, iki sütun, 9 punto ana metin.
\usepackage[utf8]{inputenc}
\usepackage[T1]{fontenc}
\usepackage[english,turkish]{babel}
\usepackage{newtxtext,newtxmath}
\usepackage{geometry}
\usepackage{titlesec}
\usepackage{caption}
\usepackage{amsmath}
\usepackage{booktabs}
\usepackage{array}
\usepackage{siunitx}
\usepackage{graphicx}
\usepackage{pgfplots}
\usepackage{tikz}
\usepackage{upgreek}
\usepackage{url}
\usepackage[hidelinks]{hyperref}
\usepackage{enumitem}
\usepackage{balance}

\pgfplotsset{compat=1.18,/pgf/number format/use comma}
\sisetup{detect-all,output-decimal-marker={,}}
\urlstyle{same}

% ELECO: ikinci ve sonraki sayfalar üst 2.5 cm; ilk sayfada başlık
% bloğuna ek 5 mm verilerek üst kenar 3 cm yapılır.
\geometry{
  a4paper,
  left=20mm,
  right=20mm,
  top=25mm,
  bottom=37mm,
  columnsep=5mm
}

\pagestyle{empty}
\setlength{\parindent}{0.4mm}
\setlength{\parskip}{0pt}
\renewcommand{\baselinestretch}{1.0}
\setlist[itemize]{nosep,leftmargin=*,topsep=0pt,partopsep=0pt}

% Başlıklar: 10 pt kalın; bölüm ortalı, alt bölüm sola hizalı.
\titleformat{\section}
  {\centering\bfseries\fontsize{10}{10.8}\selectfont}
  {\thesection.}{0.5em}{}
\titleformat{\subsection}
  {\bfseries\fontsize{10}{10.8}\selectfont}
  {\thesubsection.}{0.5em}{}
\titlespacing*{\section}{0pt}{9pt}{9pt}
\titlespacing*{\subsection}{0pt}{9pt}{9pt}

% Şekil/çizelge başlıkları: 9 pt, iki yana yaslı.
\captionsetup{
  font=normalfont,
  labelfont=bf,
  justification=justified,
  singlelinecheck=false
}
\captionsetup[figure]{name=Şekil,labelsep=colon,position=bottom,skip=6pt}
\captionsetup[table]{name=Çizelge,labelsep=period,position=top,skip=6pt}
\setlength{\textfloatsep}{9pt}
\setlength{\floatsep}{9pt}
\setlength{\intextsep}{9pt}

% thebibliography'nin kendi başlığını kaldır; başlık ELECO biçiminde
% numaralı \section{Kaynaklar} olarak ayrıca verilir.
\makeatletter
\renewenvironment{thebibliography}[1]
  {\list{\@biblabel{\@arabic\c@enumiv}}%
     {\settowidth\labelwidth{\@biblabel{#1}}%
      \leftmargin\labelwidth
      \advance\leftmargin\labelsep
      \usecounter{enumiv}%
      \let\p@enumiv\@empty
      \renewcommand\theenumiv{\@arabic\c@enumiv}}%
   \sloppy
   \clubpenalty4000
   \widowpenalty4000
   \setlength{\itemsep}{0pt}%
   \setlength{\parsep}{0pt}%
   \setlength{\parskip}{0pt}}
  {\endlist}
\makeatother

\begin{document}
\selectlanguage{turkish}
\shorthandoff{=}

\twocolumn[
\begin{center}
  \vspace*{5mm}
  {\fontsize{14}{14}\selectfont\bfseries
  Erişimi Kısıtlı Hassas Zamanlama Donanımı için Yüksek Doğruluklu Davranışsal Simülatör:\\
  Rubidyum Osilatörü Üzerinde Bir Vaka Çalışması\par}
  \vspace{2pt}
  {\fontsize{14}{14}\selectfont\bfseries
  A High-Fidelity Behavioral Simulator for Limited-Access Precision Timing Hardware:\\
  A Case Study on a Rubidium Oscillator\par}

  \vspace{9pt}
  {\fontsize{11}{11}\selectfont
  Ömer Faruk Karagöz$^{1}$, Ömer Yıldız$^{1}$, Batuhan Canlı$^{1}$\par}

  \vspace{9pt}
  {\fontsize{10}{10.8}\selectfont
  $^{1}$TÜBİTAK BİLGEM MA3, Türkiye\\
  omer.karagoz@std.bogazici.edu.tr, omer.yildiz@tubitak.gov.tr, ahmet.canli@tubitak.gov.tr\par}
  \vspace{18pt}
\end{center}
]

\begin{center}
  {\fontsize{10}{10.8}\selectfont\bfseries Özet\par}
\end{center}
\vspace{-6pt}
{\fontsize{9}{10.8}\selectfont\itshape
\noindent Hassas zamanlama donanımına bağımlı istemci yazılımlarının geliştirilmesi, cihazın pahalı, sınırlı sayıda veya belirli çalışma durumlarına ulaşmak için uzun süre gerektiren bir donanım olması nedeniyle güçleşebilir. Bu çalışmada, ticari bir düşük gürültülü rubidyum Küresel Navigasyon Sistemi Disiplinli Osilatörü (\emph{Global Navigation System Disciplined Oscillator}, GNSDO) için yüksek doğruluklu davranışsal bir simülatör sunulmaktadır. Simülatör, cihazın seri arayüzünü, Programlanabilir Cihazlar için Standart Komutlar (\emph{Standard Commands for Programmable Instruments}, SCPI) protokolünü, durum geçişlerini ve istemci tarafından gözlenebilen zamana bağlı davranışlarını yeniden üretmektedir. Model, üretici kılavuzu, sınırlı cihaz kayıtları ve doğrudan gözlemler kullanılarak oluşturulmuştur. Holdover sırasında zaman aralığı farkı, ısınma ve kilitlenme geçişleri, sağlık bayrakları ve ölçüm değişimleri modellenmiştir. Doğrulama; kaydedilmiş yanıtlarla karşılaştırma, parametre izlenebilirliği, değişkenler arası ilişkiler ve otomatik değişmez testleriyle gerçekleştirilmiştir. Sonuçlar, simülatörün gerçek donanıma sürekli erişim gerektirmeden tekrarlanabilir bir test ortamı sağlayabildiğini göstermektedir.\par}

\vspace{9pt}
\begin{center}
  {\fontsize{10}{10.8}\selectfont\bfseries Abstract\par}
\end{center}
\vspace{-6pt}
{\selectlanguage{english}\fontsize{9}{10.8}\selectfont\itshape
\noindent Developing client software that depends on precision timing hardware can be difficult when the device is expensive, available only in limited numbers, or requires long periods to reach specific operating states. This paper presents a high-fidelity behavioral simulator for a commercial low-noise rubidium Global Navigation System Disciplined Oscillator (GNSDO). The simulator reproduces the device's serial interface, Standard Commands for Programmable Instruments (SCPI) protocol, state transitions, and time-dependent behavior observable by client software. The model was constructed from the vendor manual, a limited set of device recordings, and direct observations. Time-interval difference during holdover, warm-up and lock transitions, health flags, and measurement variations are modeled. Validation is performed through comparison with recorded responses, parameter traceability, relationships among variables, and automated invariant tests. The results show that the simulator can provide a repeatable test environment without requiring continuous access to the physical device.\par}
\selectlanguage{turkish}

\section{Giriş}
Küresel Navigasyon Uydu Sistemi (\emph{Global Navigation Satellite System}, GNSS) ile disipline edilen osilatörler, telekomünikasyon, ölçüm, enerji ve savunma uygulamalarında izlenebilir zaman ve frekans üretmek için kullanılmaktadır. Bu cihazlar, GNSS sisteminin uzun dönem doğruluğunu yerel bir osilatörün kısa dönem kararlılığıyla birleştirir. İncelenen cihazda yerel referans rubidyum tabanlı bir atomik osilatördür ve cihazın izleme, yapılandırma ve durum bilgilerine seri arayüz üzerinden, Programlanabilir Cihazlar için Standart Komutlar (\emph{Standard Commands for Programmable Instruments}, SCPI) kullanılarak erişilebilmektedir \cite{jlt2015,scpi1999}.

Bu tür bir cihaza bağımlı istemci yazılımlarının geliştirilmesinde fiziksel erişim önemli bir darboğaz oluşturabilir. Donanımın pahalı veya sınırlı sayıda olması, aynı cihazın birden fazla geliştirici ya da sürekli entegrasyon ortamı tarafından eşzamanlı kullanılamaması ve GNSS kaybı, referans kaybı sonrasındaki serbest çalışma (\emph{holdover}) veya ilk açılış gibi durumların kontrollü olarak tekrar edilmesinin zaman alması test kapsamını sınırlar. Gerçek donanımın test ortamında bulunmaması veya kullanımının zor olması, yazılım testinde alternatif uygulamalar ve donanım simülatörleri kullanılmasının temel gerekçelerinden biridir \cite{meszaros2007,morita2017}.

Bu çalışmada, üreticinin \emph{Global Navigation System Disciplined Oscillator} (GNSDO) olarak adlandırdığı Jackson Labs Low Noise Rubidium cihazı için, istemci yazılımının gördüğü seri protokolü ve zamana bağlı davranışı yeniden üreten bir \emph{high-fidelity behavioral simulator} geliştirilmiştir. Buradaki ``high-fidelity'' ifadesi, simülatörün kullanım amacına göre tanımlanmıştır: hedef, cihazın faz gürültüsü veya Allan sapması gibi metrolojik özelliklerini yeniden üretmek değil; istemci yazılımının gözlemlediği komut/yanıt biçimini, tutarlı cihaz durumunu ve ilgili zaman alanı davranışlarını yeterli doğrulukta yeniden üretmektir. Simülasyon doğrulamasında modelin geçerliliğinin kullanım amacıyla birlikte ele alınması gerektiği literatürde de vurgulanmaktadır \cite{sargent2013}.

Çalışmanın temel katkıları şunlardır:
\begin{itemize}
    \item Gerçek seri aygıtla aynı erişim biçimini sağlayan sanal seri port, SCPI ayrıştırıcı, tek bir cihaz-durumu modeli ve zamana bağlı davranış katmanlarından oluşan modüler bir mimari;
    \item holdover, ısınma, kilitlenme, sağlık bayrakları ve ölçüm değişkenleri için tekrarlanabilir davranış modelleri;
    \item donanıma sınırlı erişim koşullarında uygulanabilen, kaynak izlenebilirliği, kayıt karşılaştırması, ilişki kontrolleri ve fiziksel değişmez testlerini birleştiren doğrulama yaklaşımı.
\end{itemize}

\section{Arka Plan ve İlgili Çalışmalar}
\subsection{Referans Cihaz ve Komut Arayüzü}
Referans cihaz Jackson Labs Technologies Low Noise Rubidium GNSDO'dur \cite{jlt2015}. Üretici, ürün adında GNSDO kısaltmasını kullanmaktadır; bu nedenle makalede ürün adı için GNSDO, genel teknoloji sınıfı için ise GNSS ile disipline edilen osilatör ifadesi kullanılmaktadır.
Kılavuz ayrıca GPS, GNSS, GNSDO ve GPSDO terimlerini eş anlamlı kullanmakta; komut ve ekran çıktılarında geçen CSAC ifadesinin geriye dönük uyumluluk nedeniyle rubidyum MAC (\emph{Miniature Atomic Clock}) osilatörü ile eş anlamlı olduğunu belirtmektedir \cite{jlt2015}. Bu nedenle makalede geçen CSAC etiketli komut ve ölçümler, referans cihazdaki rubidyum osilatörüne aittir. Üretici kılavuzuna göre cihazın RS-232 seri arayüzünün fabrika ayarı 115200 baud hız, 8 veri biti, parite yok, 1 stop biti ve akış kontrolü yok biçimindedir; bu yapı kısaca 8N1 olarak gösterilmektedir. Komut arayüzü SCPI sözdizimini kullanır; SCPI standardı programlanabilir ölçüm cihazları için donanım arayüzünden bağımsız bir komut dili tanımlar \cite{scpi1999}.

Cihazda rubidyum osilatörü GNSS veya harici saniyede bir darbe (\emph{one pulse per second}, 1PPS) referansına kilitleyen ve filtre osilatörünü rubidyum referansına kilitleyen iki ayrı servo döngüsü bulunmaktadır \cite{jlt2015}. \texttt{SYNC:TINT?} komutu, seçili döngü için \emph{Time Interval} (TINT) olarak adlandırılan 1PPS faz farkını raporlar. \texttt{SYNC:FEEstimate?} komutu ise üreticinin 1000~s ölçüm aralığında Allan varyansına benzer bir büyüklük olarak tanımladığı \emph{Frequency Error Estimate} (FEE; frekans hata kestirimi) değerini döndürür. \texttt{SYNC:HEALTH?} komutu çalışma süresi, holdover ve faz farkı gibi koşullara bağlı sağlık bayraklarını onaltılık bir bit alanı olarak verir. Tanısal komutlarda kullanılan \emph{Electronic Frequency Control} (EFC; elektronik frekans kontrolü) değerleri de iç referans osilatörüne uygulanan kontrol büyüklüğünü ifade eder \cite{jlt2015}.

\subsection{Davranışsal Simülasyon ve Test Amaçlı Donanım İkamesi}
Yazılım testi literatüründe \emph{test double}, gerçek bir bağımlılığın test ortamında kullanılamadığı, çok yavaş olduğu veya test maliyetini artırdığı durumlarda onun yerine kullanılan test bileşenlerini tanımlayan genel bir terimdir. Meszaros'un sınıflandırmasında, gerçek bileşenin işlevini daha basit fakat çalışan bir uygulamayla sağlayan test ikamesi \emph{Fake Object} olarak adlandırılır \cite{meszaros2007}. Bu çalışmanın teknik odağı test ikamesi taksonomisinden çok davranışsal simülasyon olduğundan, devamında ``simülatör'' terimi kullanılmaktadır.

Üretim gözlemlerinden gerçekçi test ikameleri üretme fikri de yazılım testi literatüründe incelenmiştir. Tiwari ve arkadaşları, üretim etkileşimlerini gözleyerek gerçekçi mock ve stub davranışları üretmenin uygulanabilirliğini göstermiştir \cite{tiwari2024}. Donanım tarafında Morita ve arkadaşları, gerçek makine olmadan yazılım doğrulaması yapabilmek için donanım ve hata davranışlarını simüle eden bir test ortamı geliştirmiştir \cite{morita2017}. Bu çalışmalar, gerçek donanımın erişilebilirliğinin sınırlı olduğu durumlarda gözlenebilir davranışın ikame edilmesinin yazılım geliştirme açısından pratik değerini desteklemektedir.

\subsection{Simülasyon Doğrulaması}
Bir simülasyon modelinin yalnızca çalışması, modelin kullanım amacı için yeterince doğru olduğunu göstermez. Sargent; kavramsal model geçerliliği, uygulama doğrulaması, veri geçerliliği ve operasyonel geçerliliği ayrı değerlendirme alanları olarak ele almakta ve modelin amaçlanan kullanım bağlamında doğrulanmasını önermektedir \cite{sargent2013}. Bu çalışmada doğrudan zaman hizalı ve kapsamlı bir donanım karşılaştırması bulunmadığından, doğrulama tek bir metriğe indirgenmemiş; birbirini tamamlayan kanıt türleri birlikte kullanılmıştır.

\subsection{Saat ve Holdover Modellemesi}
Saat fazı ve frekans hatası arasındaki ilişki klasik saat modellerinin temelidir. Zucca ve Tavella, saat hatasını stokastik diferansiyel denklemlerle ifade ederek frekans süreçleri ile Allan varyansı arasındaki ilişkiyi ayrıntılı biçimde incelemiştir \cite{zucca2005}. Bu çalışmanın amacı osilatörün tam stokastik spektrumunu modellemek değildir. Bunun yerine, istemci testleri için gerekli zaman alanı davranışları daha düşük karmaşıklıkta ve deterministik biçimde temsil edilmektedir.

\section{Tasarım Hedefleri ve Kapsam}
\label{sec:goals}
Simülatörün temel hedefi, gerçek cihazla çalışan istemci yazılımında değişiklik gerektirmeden aynı bağlantı ve komut modelini sunmaktır. Bu amaçla aşağıdaki doğruluk boyutları esas alınmıştır:
\begin{enumerate}
    \item \textbf{Arayüz doğruluğu:} Komut adları, kısa/uzun SCPI biçimleri, satır sonlandırma ve çok satırlı yanıtların biçimi;
    \item \textbf{Durum tutarlılığı:} Özet sorgular ile tekil sorguların aynı cihaz durumundan türetilmesi;
    \item \textbf{Zamansal doğruluk:} Isınma, kilitlenme, holdover ve sağlık geçişlerinin gerçek cihazla aynı mertebede ilerlemesi;
    \item \textbf{İlişkisel doğruluk:} Birbirine bağlı ölçümlerin kayıtlarda ve kılavuzda gözlenen ilişkileri koruması;
    \item \textbf{Tekrarlanabilirlik:} Aynı başlangıç koşulları altında aynı test çalışmasının aynı çıktıları üretmesi.
\end{enumerate}

Spektral faz gürültüsü, Allan sapması eğrisinin tam yeniden üretimi ve cihazlar arası üretim varyasyonu bu çalışmanın kapsamı dışındadır.

\section{Simülatör Mimarisi}
Simülatör dört katmandan oluşmaktadır:
\begin{itemize}
    \item \textbf{Sanal seri arayüz:} \texttt{socat} ile oluşturulan sözde terminal (\emph{pseudo-terminal}) üzerinden istemciye standart bir \texttt{/dev/tty*} aygıtı sunulur.
    \item \textbf{SCPI ayrıştırıcı:} Komutlar kısa ve uzun biçimleriyle normalize edilir, uygun işleyiciye yönlendirilir ve desteklenmeyen girişlerde cihazla uyumlu hata yanıtı üretilir.
    \item \textbf{Cihaz-durumu katmanı:} Kimlik, yapılandırma, servo durumu, ölçümler ve sağlık bilgileri tek bir durum nesnesinde tutulur. Böylece birbiriyle ilişkili sorgular aynı anda çelişkili değer üretmez.
    \item \textbf{Davranış modelleri:} Zamana bağlı fiziksel ve mantıksal davranışlar giriş/çıkış katmanından ayrılmış saf fonksiyonlar olarak uygulanır.
\end{itemize}

Davranış modellerinin geçen süreyi doğrudan parametre olarak alması, saatler süren gerçek zamanı beklemeden uzun holdover veya ilk açılış senaryolarının birim testinde hızlandırılmış biçimde sınanabilmesini sağlar.

\section{Davranış Modelleri}
\subsection{Holdover Sırasında Zaman Aralığı Farkı}
GNSS referansı kaybedildiğinde servo döngüsü referansa göre düzeltme yapamaz ve yerel osilatörün artık frekans hatası faz hatasına dönüşür. Kısa ve orta süreli istemci testleri için birinci dereceden yaklaşım kullanılmıştır:
\begin{equation}
    x(t)=x_0+y_h t,
    \label{eq:holdover}
\end{equation}
burada $x(t)$ zaman aralığı farkı, $x_0$ holdover başlangıcındaki faz farkı ve $y_h$ etkin artık fraksiyonel frekans hatasıdır.

Üretici kılavuzu \texttt{SYNC:FEEstimate?} değerini 1000~s aralığında Allan varyansına benzer bir büyüklük olarak tanımladığından, bu değer doğrudan işaretli fiziksel frekans sapması olarak yorumlanmamıştır. Simülatörde $|y_h|$ için kaydedilmiş FEE büyüklükleri başlangıç referansı olarak kullanılmakta; büyüklük ve işaret, senaryo yapılandırmasından ayarlanabilen açık bir model parametresi olarak tutulmaktadır. Bu yaklaşım, kaynak tarafından desteklenmeyen bir fiziksel yorumu modele gömmek yerine varsayımı görünür kılar.
Kılavuza göre \texttt{SYNC:HOLD:INIT} ile başlatılan zorlanmış holdover sırasında GNSS sinyali alınmaya devam ettiğinden cihaz, 1PPS çıkışı ile GNSS 1PPS sinyali arasındaki zaman aralığı farkını raporlamayı sürdürür \cite{jlt2015}. Denklem~\ref{eq:holdover} bu gözlenebilir büyüklüğe karşılık gelir; referans sinyalinin fiziksel olarak kaybolduğu durumda raporlanan değerler kayıtlarla doğrulanmamıştır.

Önceki sürümde kullanılan ikinci dereceden yaşlanma terimi bu çalışmadan çıkarılmıştır. Üretici kılavuzundaki $8\times10^{-14}$ değeri yaşlanma katsayısı değil, belirli bir ortalama süresindeki Allan sapması performansını ifade etmektedir; bu nedenle bu değerin $D$ yaşlanma katsayısı olarak kullanılması izlenebilir değildir.

\subsection{Deterministik Ölçüm Değişimi}
Gerçek cihazda sıcaklık, TINT ve benzeri büyüklükler tamamen sabit değildir. Ancak her sorguda bağımsız rastgele sayı üretmek hem sorgu sıklığına bağlı yapay değişim oluşturur hem de regresyon testlerinin tekrar üretilebilirliğini bozar. Bu nedenle gözlenebilir küçük değişimler zamanın deterministik bir fonksiyonu olarak modellenmiştir:
\begin{equation}
 v(t)=v_0+A\frac{\sum_{k=0}^{2}w_k\sin\!\left(2\pi r_k t/P+\upvarphi_k\right)}{\sum_{k=0}^{2}w_k},
 \label{eq:variation}
\end{equation}
burada $v_0$ nominal değer, $A$ izin verilen maksimum sapma, $w_k\in\{1,1/2,1/4\}$ ağırlıklar, $r_k\in\{1,\upphi,\upphi^2\}$ birbirine tam kat olmayan frekans oranları ve $\upvarphi_k$ kanal bazlı sabit faz ofsetleridir. Normalizasyon sayesinde değişim $\pm A$ aralığıyla sınırlı kalır. Bu yapı fiziksel gürültünün spektral modeli olarak değil, istemci testleri için zaman-tutarlı ve tekrarlanabilir küçük değişim üretmek amacıyla kullanılmaktadır.

\subsection{Isınma ve Kilitlenme}
Baskılı devre kartı (\emph{printed circuit board}, PCB) sıcaklığı için birinci dereceden termal yaklaşım kullanılmıştır:
\begin{equation}
 T(t)=T_{\infty}-\left(T_{\infty}-T_0\right)e^{-t/\uptau},
 \label{eq:thermal}
\end{equation}
burada $T_0$ başlangıç sıcaklığı, $T_{\infty}$ kararlı durum sıcaklığı ve $\uptau$ deneysel termal zaman sabitidir. Varsayılan parametreler, eldeki gözlem ve kayıtlarla uyumlu olacak şekilde $T_0=\SI{25}{\celsius}$, $T_{\infty}\approx\SI{52.8}{\celsius}$ ve $\uptau\approx\SI{90}{s}$ olarak seçilmiştir. $T_0$ ortam koşuluna bağlı olduğu için yapılandırılabilir tutulmuştur; $\uptau$ üretici spesifikasyonu değil, gözlemsel bir model parametresidir.

Üretici kılavuzu, GNSS anteni bağlıyken tam kilit göstergesinin güç verildikten sonra tipik olarak 20 dakikadan kısa sürede oluştuğunu belirtmektedir \cite{jlt2015}. Servo durumları bu zaman ölçeğini koruyacak şekilde ısınma, kilitlenme ve kilitli durumlara ayrılmıştır. Kılavuzdaki ``2 dakikadan kısa sürede atomik kilit'' ifadesi, çip ölçekli atomik saat (\emph{Chip-Scale Atomic Clock}, CSAC) kullanılan belirli varyanta ait olduğundan referans rubidyum cihazının genel ısınma süresi olarak kullanılmamıştır.

\subsection{Sağlık Durumu}
\texttt{SYNC:HEALTH?} için sağlık bayrakları üretici kılavuzundaki koşullara göre oluşturulmaktadır. Modelde özellikle aşağıdaki kurallar uygulanmaktadır: çalışma süresi 200~s'den küçükse \texttt{0x8}; holdover süresi 60~s'yi aşarsa \texttt{0x10}; referansa göre faz farkı 210~ns'yi aşarsa \texttt{0x4}. Bayraklar bit düzeyinde mantıksal VEYA işlemiyle birleştirilir. Örneğin uzun holdover sırasında faz farkının 210~ns'yi aşması \texttt{0x14} sonucunu üretir. Sağlık bayrağındaki 210~ns koşulu ile \texttt{SYNC:TINT:THR} komutunun 220~ns varsayılan ani faz hizalama (\emph{jam-sync}) eşiği ayrı kavramlardır ve modelde birbirine karıştırılmamıştır \cite{jlt2015}.

\section{Modelin İnşası ve Parametre İzlenebilirliği}
Model parametreleri dört kaynak sınıfından biriyle ilişkilendirilmiştir: üretici kılavuzu, kaydedilmiş cihaz çıktısı, doğrudan operatör gözlemi ve açık varsayım. Böylece her sayısal değerin kaynağı ve güven düzeyi incelenebilir hale gelir. Çizelge~\ref{tab:trace} bu ayrımı özetlemektedir.

\begin{table}[t]
\caption{Model parametrelerinin izlenebilirliği}
\label{tab:trace}
\centering
\fontsize{9}{10.8}\selectfont
\setlength{\tabcolsep}{3pt}
\begin{tabular}{>{\raggedright\arraybackslash}p{0.33\columnwidth}>{\raggedright\arraybackslash}p{0.25\columnwidth}>{\raggedright\arraybackslash}p{0.33\columnwidth}}
\toprule
\textbf{Parametre / ilişki} & \textbf{Değer} & \textbf{Kaynak ve durum} \\
\midrule
Seri arayüz & 115200, 8N1 & Kılavuz, doğrudan \\
Sağlık faz koşulu & $|\mathrm{TINT}|>210$ ns & Kılavuz, doğrudan \\
Ani faz hizalama (jam-sync) eşiği & 220 ns varsayılan & Kılavuz, doğrudan \\
Holdover bayrağı & $t>60$ s & Kılavuz, doğrudan \\
İlk çalışma bayrağı & $t<200$ s & Kılavuz, doğrudan \\
$|y_h|$ başlangıç aralığı & $1{,}3$--$2{,}5\times10^{-11}$ & Kayıt + model eşleme varsayımı \\
Kilitli TINT değişimi & yaklaşık $\pm12$ ns & Kayıt \\
CSAC--PCB sıcaklık farkı & $+1{,}32\,^{\circ}$C & Kayıtlar (4 çift) \\
EFC mutlak/bağıl ilişki & $\times200$ & Kayıtlar \\
$\uptau$ & yaklaşık 90 s & Operatör gözlemi / model parametresi \\
$T_0$ & $25\,^{\circ}$C & Varsayım, yapılandırılabilir \\
\bottomrule
\end{tabular}
\end{table}

Bu tablo aynı zamanda modelin kanıt düzeyini sınırlayan noktaları açıkça göstermektedir. Özellikle FEE değerinden etkin holdover eğimine yapılan eşleme ve termal zaman sabiti, doğrudan üretici spesifikasyonu olarak sunulmamaktadır.

Kılavuz ve kayıtların birlikte kullanılması, yalnızca kılavuza dayanan bir modelde gözden kaçabilecek noktaları da ortaya çıkarmıştır (Çizelge~\ref{tab:findings}). Uydu sayılarındaki durum ise bilinçli olarak taklit edilmemiş, istemci testlerinde tutarlı bir durum sunmak için modelde sınırlandırılmıştır.

\begin{table}[t]
\caption{Kılavuz ve kayıtların birlikte değerlendirilmesinden çıkan bulgular}
\label{tab:findings}
\centering
\fontsize{9}{10.8}\selectfont
\setlength{\tabcolsep}{3pt}
\begin{tabular}{>{\raggedright\arraybackslash}p{0.22\columnwidth}>{\raggedright\arraybackslash}p{0.30\columnwidth}>{\raggedright\arraybackslash}p{0.39\columnwidth}}
\toprule
\textbf{Konu} & \textbf{Kılavuz} & \textbf{Kayıt ve model} \\
\midrule
TINT & $<0{,}2$ ns ortalama faz doğruluğu & İşaretli, yaklaşık $\pm12$ ns; 0,2 ns ortalama kabul edildi \\
\texttt{MEAS:CURR?} & Eski komut; akım yerine sıcaklık & Kayıtta sıcaklık; model de sıcaklık döndürür \\
EFC & Bağıl (\%) ve mutlak (ppt) & Mutlak $=200\times$ bağıl (4/4) \\
Uydu sayısı & İzlenen ve görünür ayrı & İzlenen (20) $>$ görünür (19); modelde sınırlandı \\
\bottomrule
\end{tabular}
\end{table}

\section{Doğrulama Yöntemi}
Doğrulama, kullanım amacına göre dört tamamlayıcı yöntemle yürütülmüştür.

\subsection{Komut ve Yanıt Karşılaştırması}
Kaydedilmiş gerçek cihaz çıktıları ile simülatör yanıtları aynı komutlar için karşılaştırılmıştır. Kimlik, seri numarası ve sabit yapılandırma alanlarında tam metin eşleşmesi; çok satırlı komutlarda alan sırası, etiketler ve satır yapısı kontrol edilmiştir.
Karşılaştırmada kullanılan kayıtlar, oturumun büyük bölümünde cihazın GNSS'e kilitli olduğu ve alıcının konum ölçümünün (\emph{survey}) sürdüğü yaklaşık 40 dakikalık tek bir oturumda alınmıştır; kayıtların bir kısmı zaman damgası içermemektedir.

\subsection{İlişki Tutarlılığı}
Tek tek değerlerin yanı sıra değişkenler arası ilişkiler doğrulanmıştır. Örneğin kaydedilmiş dört ölçümde EFC mutlak değeri bağıl değerin 200 katı olarak gözlenmiş, CSAC sıcaklığı ise PCB sıcaklığından yaklaşık \SI{1.32}{\celsius} yüksek ölçülmüştür. Simülatör bu ilişkileri ayrı sabit yanıtlar yerine ortak durum üzerinden üretmektedir.

\subsection{Fiziksel ve Mantıksal Değişmezler}
Model için otomatik test edilen başlıca değişmezler; sıcaklığın kararlı değere sınırlı yaklaşması, aynı başlangıç koşullarında aynı serinin üretilmesi, holdover süresinin monoton artması, sağlık bayraklarının koşullarla tutarlı olması ve özet/tekil sorguların aynı cihaz durumunu yansıtmasıdır.

\subsection{Zamansal Davranış Kontrolü}
Uzun süreli durumlar gerçek zaman beklemeden hızlandırılmış olarak taranmıştır. Böylece 20 dakikalık kilitlenme zaman ölçeği, saatler süren holdover ve sağlık eşiği geçişleri birim testleri içinde değerlendirilebilmektedir.

\section{Sonuçlar}
\subsection{Arayüz ve İlişki Doğruluğu}
Çizelge~\ref{tab:compare}, mevcut kayıt kümesi üzerinde yapılan karşılaştırmaların bir bölümünü göstermektedir.

\begin{table}[t]
\caption{Simülatör çıktısının kaydedilmiş cihaz çıktısıyla karşılaştırılması}
\label{tab:compare}
\centering
\fontsize{9}{10.8}\selectfont
\setlength{\tabcolsep}{3pt}
\begin{tabular}{p{0.26\columnwidth}p{0.44\columnwidth}p{0.20\columnwidth}}
\toprule
\textbf{Sorgu} & \textbf{Kayıt / beklenen ilişki} & \textbf{Sonuç} \\
\midrule
\texttt{*IDN?} & Jackson Labs, LN Rb GPSDO (PRE), Firmware Rev 1.17 & Tam \\
\texttt{SYST:ID:SN?} & 122301668 & Tam \\
\texttt{CSAC:SN?} & 2209MX04906 & Tam \\
\texttt{CSAC:FW?} & V1.0.23 & Tam \\
\texttt{MEAS?} & dört etiketli satır & Biçim tam \\
\texttt{DIAG?} & $\mathrm{EFC}_{\mathrm{abs}}=200\,\mathrm{EFC}_{\mathrm{rel}}$ & 4/4 \\
CSAC--PCB & $+1{,}32\,^{\circ}$C & $\pm0{,}07\,^{\circ}$C \\
\texttt{SYNC:FEE?} & $1{,}3$--$2{,}5\times10^{-11}$ & Aralık \\
\texttt{SYNC:TINT?} & işaretli, yaklaşık $\pm12$ ns & Aralık/işaret \\
\bottomrule
\end{tabular}
\end{table}

Kimlik ve yapılandırma alanlarında tam eşleşme elde edilmiştir. Türetilmiş ölçümlerde ise sabit değer kopyalamak yerine gözlenen ilişkilerin korunması tercih edilmiştir. Bu yaklaşım, simülatörün farklı komutlara verdiği yanıtların aynı iç durumla tutarlı kalmasını sağlamaktadır.

\subsection{Holdover Davranışı}
Örnek bir $|y_h|=1{,}8\times10^{-11}$ değeri için Denklem~\ref{eq:holdover}, bir saat sonunda yaklaşık 64,8~ns ve bir gün sonunda yaklaşık \SI{1.56}{\micro\second} zaman aralığı farkı üretmektedir. 210~ns sağlık koşulu yaklaşık 3,24~saatte aşılmaktadır. Kaydedilen FEE aralığı ($1{,}3$--$2{,}5\times10^{-11}$) kullanıldığında bu süre yaklaşık 2,3--4,5~saat arasında değişmektedir. Şekil~\ref{fig:holdover} bu davranışı göstermektedir.

\begin{figure}[t]
\centering
\begin{tikzpicture}
\begin{axis}[
 width=\columnwidth,
 height=0.62\columnwidth,
 xlabel={Holdover süresi (saat)},
 ylabel={TINT (ns)},
 xmin=0,xmax=12,
 ymin=0,ymax=800,
 grid=both,
 tick label style={font=\scriptsize},
 label style={font=\small},
]
\addplot[thick,domain=0:12,samples=2] {64.8*x};
\addplot[dashed,domain=0:12,samples=2] {210};
\addplot[only marks,mark=*] coordinates {(3.24074,210)};
\node[font=\scriptsize,anchor=north west] at (axis cs:3.4,195) {3,24 saat};
\node[font=\scriptsize,anchor=south east] at (axis cs:11.8,210) {210 ns sağlık koşulu};
\end{axis}
\end{tikzpicture}
\caption{$|y_h|=1{,}8\times10^{-11}$ için doğrusal holdover zaman aralığı farkı}
\label{fig:holdover}
\end{figure}

Bu hesap bir metrolojik holdover performans garantisi değildir; yalnızca simülatörde kullanılan etkin frekans hata parametresinin istemci tarafından gözlenen TINT değerine nasıl yansıtıldığını gösterir.

\subsection{Sağlık Durumu Geçişleri}
Holdover 60~s'yi aştığında \texttt{0x10} bayrağı etkinleşir. Faz farkı 210~ns'yi aştığında \texttt{0x4} eklenir ve birleşik durum \texttt{0x14} olur. Güç verildikten sonraki ilk 200~s için \texttt{0x8} bayrağı ayrıca modellenmektedir. Bu kurallar üretici kılavuzundaki sağlık bitleriyle doğrudan uyumludur \cite{jlt2015}.

\subsection{Isınma ve Kilitlenme}
Denklem~\ref{eq:thermal}'teki varsayılan parametrelerle PCB sıcaklığı birkaç dakika içinde kararlı değere yaklaşırken, GNSS kilit durumu daha yavaş ilerlemekte ve modelde yaklaşık 20 dakikalık zaman ölçeğinde etkinleşmektedir. Şekil~\ref{fig:warmup} sıcaklık modeli ile kilit geçişinin farklı zaman ölçeklerini göstermektedir.

\begin{figure}[t]
\centering
\begin{tikzpicture}
\begin{axis}[
 width=\columnwidth,
 height=0.62\columnwidth,
 xlabel={Açılıştan itibaren süre (dk)},
 ylabel={PCB sıcaklığı ($^{\circ}$C)},
 xmin=0,xmax=25,
 ymin=24,ymax=55,
 grid=both,
 tick label style={font=\scriptsize},
 label style={font=\small},
]
\addplot[thick,domain=0:25,samples=120] {52.8-(52.8-25)*exp(-x*60/90)};
\addplot[dashed] coordinates {(20,24) (20,55)};
\node[font=\scriptsize,anchor=south,rotate=90] at (axis cs:19.4,38) {yaklaşık GNSS kilidi};
\end{axis}
\end{tikzpicture}
\caption{Birinci dereceden PCB sıcaklık modeli ve yaklaşık GNSS kilit zaman ölçeği}
\label{fig:warmup}
\end{figure}

\subsection{Doğruluk Boyutlarına Göre Özet}
Çizelge~\ref{tab:dims}, Bölüm~\ref{sec:goals}'te tanımlanan doğruluk boyutlarının hangi kanıtla ve hangi sınırlarla karşılandığını özetlemektedir. Model fonksiyonları geçen süre parametre olarak verilerek tekrar çalıştırıldığında aynı başlangıç koşulları için üretilen serilerin maksimum farkı sıfırdır. Seri arayüz üzerinden yürütülen testlerde ise yanıtlar sorgunun geliş anına bağlı olduğundan, tam tekrar için test zamanının kontrol edilmesi gerekir.

\begin{table}[t]
\caption{Doğruluk boyutlarına göre doğrulama durumu}
\label{tab:dims}
\centering
\fontsize{9}{10.8}\selectfont
\setlength{\tabcolsep}{3pt}
\begin{tabular}{>{\raggedright\arraybackslash}p{0.18\columnwidth}>{\raggedright\arraybackslash}p{0.30\columnwidth}>{\raggedright\arraybackslash}p{0.43\columnwidth}}
\toprule
\textbf{Boyut} & \textbf{Kanıt} & \textbf{Durum ve sınır} \\
\midrule
Arayüz & Kayıt; 66 \texttt{HELP?} sorgusu & Kayıttaki biçimler eşleşti; 50/66 sorgu uygulandı \\
Durum & Özet/tekil sorgu testleri & Aynı cihaz durumundan türetiliyor \\
Zamansal & Kılavuz zaman ölçekleri & Uyumlu; cihaz zaman serisiyle karşılaştırılmadı \\
İlişkisel & EFC, CSAC--PCB & Uyumlu; aynı kayıtlardan türetildi \\
Tekrar & Model çalıştırmaları & Model: fark 0; seri arayüz: sorgu anına bağlı \\
\bottomrule
\end{tabular}
\end{table}

\section{Sınırlamalar}
Simülatörün yüksek doğruluk iddiası istemci yazılımının gözlemlediği arayüz, durum ilişkileri ve zaman alanı davranışlarıyla sınırlıdır. Faz gürültüsü, Allan sapması ve osilatörün ayrıntılı stokastik karakteri modellenmemektedir. Ölçüm değişimleri test tekrar edilebilirliği için bilinçli olarak deterministiktir; bu nedenle simülatör osilatör kararlılığının metrolojik değerlendirilmesinde kullanılamaz. Değişimler $\pm A$ ile sınırlı olduğundan, gerçek cihazda seyrek görülebilen büyük sapmalar ve bunlara bağlı eşik aşımları üretilmez. Holdover'dan dönüş ve ani faz hizalama (\emph{jam-sync}) davranışı da modellenmemiştir.

Doğrulama tek bir fiziksel cihaz örneği ve sınırlı kayıt kümesiyle yapılmıştır. Bazı kayıtlarda zaman damgası bulunmadığından tüm davranışlar zaman hizalı gerçek/simülatör eğrileriyle karşılaştırılamamıştır. Ayrıca FEE değerinden etkin holdover eğimine yapılan eşleme üretici kılavuzunun doğrudan tanımladığı bir fiziksel ilişki değildir ve model parametresi olarak açıkça tutulmaktadır. Termal zaman sabiti de üretici spesifikasyonu değil, gözleme dayalı bir parametredir.

\balance
\section{Sonuç ve Gelecek Çalışmalar}
Bu çalışmada erişimi kısıtlı bir hassas zamanlama cihazı için seri/SCPI arayüzünü, tutarlı cihaz durumunu ve istemci açısından önemli zamana bağlı davranışları yeniden üreten yüksek doğruluklu davranışsal bir simülatör geliştirilmiştir. Simülatör, gerçek donanımın sürekli olarak test ortamında bulunmasına gerek kalmadan istemci yazılımının aynı komut arayüzü üzerinden çalıştırılmasına ve uzun süreli durumların tekrarlanabilir biçimde sınanmasına olanak vermektedir.

Çalışmanın temel sonucu, sınırlı donanım erişiminde doğruluğun tek bir ``gerçeğe benzerlik'' iddiasıyla değil; parametre kaynaklarının izlenebilirliği, kaydedilmiş yanıtlarla karşılaştırma, değişkenler arası ilişkilerin korunması ve otomatik değişmez testlerinin birlikte kullanılmasıyla daha savunulabilir hale getirilebilmesidir. Gelecek çalışmalarda zaman damgalı ve daha geniş bir cihaz kayıt kümesi toplanması, birden fazla fiziksel örnek üzerinde cihazlar arası varyasyonun ölçülmesi, model parametrelerinin otomatik kestirimi ve simülatörün sürekli entegrasyon ve sürekli teslim süreçlerindeki regresyon testlerine sistematik olarak dahil edilmesi planlanmaktadır.

\section{Erişilebilirlik}
Simülatör, ölçüm/çizim ardışık düzeni ve doğrulamada kullanılan kaydedilmiş cihaz çıktıları \url{https://github.com/Pangea567/gnsdo-scpi-simulator/tree/develop} adresinde sunulmaktadır.

\section{Kaynaklar}
\begin{thebibliography}{99}

\bibitem{jlt2015}
Jackson Labs Technologies, Inc., ``Low Noise Rubidium GNSDO User Manual,'' Document 80200526, Version 1.3, Aug. 2015.

\bibitem{scpi1999}
SCPI Consortium, \emph{1999 SCPI Syntax \& Style}. SCPI Consortium / IVI Foundation, 1999.

\bibitem{meszaros2007}
G. Meszaros, \emph{xUnit Test Patterns: Refactoring Test Code}. Addison-Wesley Professional, 2007.

\bibitem{morita2017}
K. Morita, M. Sato, and S. Wakui, ``Development of the Support System for a Software Testing Using a Hardware Failure Simulation,'' \emph{Journal of the Japan Society for Precision Engineering}, vol. 83, no. 1, pp. 101--107, 2017, doi: 10.2493/jjspe.83.101.

\bibitem{sargent2013}
R. G. Sargent, ``Verification and Validation of Simulation Models,'' \emph{Journal of Simulation}, vol. 7, no. 1, pp. 12--24, 2013, doi: 10.1057/jos.2012.20.

\bibitem{tiwari2024}
D. Tiwari, M. Monperrus, and B. Baudry, ``Mimicking Production Behavior With Generated Mocks,'' \emph{IEEE Transactions on Software Engineering}, vol. 50, no. 11, pp. 2921--2946, 2024, doi: 10.1109/TSE.2024.3458448.

\bibitem{zucca2005}
C. Zucca and P. Tavella, ``The Clock Model and Its Relationship with the Allan and Related Variances,'' \emph{IEEE Transactions on Ultrasonics, Ferroelectrics, and Frequency Control}, vol. 52, no. 2, pp. 289--296, 2005, doi: 10.1109/TUFFC.2005.1406554.

\end{thebibliography}

\end{document}
